\section{An explicit nonzero axis subspace in every characteristic}

\subsection*{C44-A04. Axes containing fibres of an abelian product}
\begin{theorem}[C44-A04]\label{thm:product-axis}
Let $r=m-1\ge2$, let $S$ and $T$ be abelian varieties of dimensions
$r$ and $r+2$, and put $X=S\times T$. Choose very ample classes
$s$ and $t$ on $S$ and $T$ and let $h=s+t$ on $X$. Let
$B\subset X$ be a smooth complete intersection of two divisors of
class $h$. Suppose that it contains the distinct fibres
$S_i=S\times\{t_i\}$ for $1\le i\le N$. Put
\[
 d_S=\int_Ss^r,\qquad d_T=\int_Tt^{r+2},\qquad
 a=\frac1{d_T}\sum_{j=0}^r(-1)^j(j+1)s^jt^{r-j},
 \qquad w_i=[S_i]_B-\kappa^*a.
 \tag{C44.A04.1}
\]
Then $w_i\in A_B$ and the exact signed Gram matrix is
\[
 \left((-1)^r\int_B w_iw_j\right)_{i,j}
       =(r+1)d_S\left(I_N-\frac1{d_T}{\bf1}{\bf1}^{t}\right).
 \tag{C44.A04.2}
\]
Smoothness of $B$ forces $N\le d_T$. Hence this matrix is positive
semidefinite in every characteristic, without 032. If $N<d_T$ it is
positive definite, and the $w_i$ are linearly independent rational
algebraic classes. In particular a single contained fibre supplies a
nonzero axis class of the required sign.

Smooth such axes containing a prescribed fibre exist after replacing
$s,t$ by sufficiently divisible positive multiples. The two defining
sections may be chosen to give a genuine Lefschetz pencil. Thus the
construction includes products of supersingular elliptic curves and
gives special axes with $A_B\ne0$ over $\overline{\F}_p$ for every $p$.
The assertion concerns the displayed subspace, not all of $A_B$.
\end{theorem}
\begin{proof}
All computations are in rational Weil cohomology, with their usual
cycle representatives. Since $s^{r+1}=0$ and $t^{r+3}=0$, the finite
identity $(1+z)^2\sum_{j=0}^r(-1)^j(j+1)z^j\equiv1\pmod{z^{r+1}}$
gives
\[
                      h^2a=t^{r+2}/d_T=[S_i]_X.
\]
Thus $\kappa_*w_i=0$ by self-intersection, and $w_i\in A_B$.
The normal bundle of $S_i$ in $X$ is trivial of rank $r+2$.
Because $B$ is smooth along $S_i$, its normal sequence is exact:
\[
 0\longrightarrow N_{S_i/B}\longrightarrow\mathcal O_S^{r+2}
   \longrightarrow L_S\oplus L_S\longrightarrow0.
\]
Here $c_1(L_S)=s$. Consequently
$c(N_{S_i/B})=(1+s)^{-2}$ and
\[
                 \int_B[S_i]^2=(-1)^r(r+1)d_S.
\]
Distinct fibres are disjoint, so their intersection is zero.
Restriction of $a$ to any $S_i$ has degree
$(-1)^r(r+1)d_S/d_T$; the same is true of
$\int_B(\kappa^*a)^2=\int_Xa h^2a$.
Expanding each $w_iw_j$ therefore gives (C44.A04.2).

We prove the bound on the number of contained fibres without using a
positivity conjecture. Write the two defining sections in a basis of
$H^0(S,L_S)$ as
$F_\nu=\sum_\alpha\sigma_\alpha\tau_{\nu\alpha}$,
where $\tau_{\nu\alpha}\in H^0(T,L_T)$ and $\nu=1,2$.
Let $Z\subset T$ be the common zero scheme of all coefficients
$\tau_{\nu\alpha}$. Its geometric points are exactly the fibres
contained in $B$. Suppose $T_zZ$ has positive dimension at a point
$z\in Z$. Differentiating the coefficients shows that the normal
surjection along $S\times\{z\}$ factors through
\[
       \mathcal O_S^{r+1}\longrightarrow L_S\oplus L_S
\]
(or through a trivial bundle of still smaller rank). Its kernel
would have rank at most $r-1$ but total Chern class $(1+s)^{-2}$.
Its $r$th Chern class would therefore be both zero, by rank, and
$(-1)^r(r+1)s^r$, whose degree is nonzero. This is impossible.
Thus $Z$ has zero tangent space at every geometric point; it is a
finite reduced scheme. Choose $r+2$ general linear combinations of
the coefficient sections. Since their common base scheme is finite,
successive dimension cuts yield a zero-dimensional complete
intersection containing $Z$. Its degree is $d_T$, so
$N\le\#Z\le d_T$.
The eigenvalues of the matrix in parentheses in (C44.A04.2) are
$1$ with multiplicity $N-1$ and $1-N/d_T$ with multiplicity one.
This proves all the asserted signs and, for $N<d_T$, independence.

For existence take the prescribed smooth centre $S_0=S\times\{0\}$.
For sufficiently positive $L_S,L_T$, the linear system
$H^0(X,I_{S_0}\otimes(L_S\boxtimes L_T))$ separates second jets
at each point off $S_0$, separates the first jets at each pair of
distinct points off $S_0$, and maps onto the conormal first jets
along $S_0$. These are uniform assertions for all such points:
apply relative Serre vanishing to the universal length-bounded
jet subschemes and the fixed coherent ideal $I_{S_0}$ (or use
high powers of product sections). A common sufficiently high
power satisfies all three conditions.
The latter assertion also follows directly from product sections:
sections of $L_S$ generate its fibres and sections of $I_0L_T$
generate the cotangent space of $T$ at $0$ after taking a high power.
For two general sections, the locus on $S_0$ where their normal
derivatives have rank at most one has expected codimension $r+1$;
the evaluation to the two normal covectors is surjective at each
point, so the incidence dimension estimate, with $\dim S_0=r$,
makes this locus empty. Away from $S_0$, jet separation and
Bertini give a smooth codimension-two intersection. Here are the
dimension counts for the pencil condition. Write $n=\dim X=2r+2$.
At a specified point off $S_0$, singularity imposes $n+1$
independent first-jet equations on a section. Degeneracy of its
quadratic part imposes an additional proper equation, by the
second-jet surjectivity. Letting the point vary gives a bad locus
of codimension at least two in the projective linear system.
Singularities at two specified distinct points impose $2(n+1)$
independent equations; their parameter space has dimension $2n$,
again giving codimension at least two. For singularity along the
imposed centre, the normal covector must vanish; this imposes
$r+2$ conditions at a specified point of $S_0$, and $\dim S_0=r$.
Its incidence image has codimension at least two as well. A general
projective line avoids these bad loci. Smoothness of its axis
also prevents the line from being tangent to the discriminant at
a single ordinary quadratic point: such tangency says that the
other pencil section vanishes at that point, putting a singular
member's zero differential into the axis's defining pair.
Thus the line gives a Lefschetz pencil, with distinct singular
values and no singularity on $S_0$. Since $n$ is even, the
nondegenerate quadratic forms are a nonempty open also in
characteristic two. These are the ordinary Lefschetz local models.
All conditions used are nonempty Zariski-open conditions over the
algebraically closed field. Over $\overline{\F}_p$ the chosen
sections are defined over a finite extension. This proves existence.
\end{proof}

\paragraph{C44-A04a: a concrete sixfold normalization.}
Take $X=E^6$, $S=E^2$, $T=E^4$ and
$h=d\sum_{i=1}^{6}[0]_i$, with $d$ sufficiently large for the
construction. For an axis containing $E^2\times\{0\}$,
\[
 d_S=2d^2,\qquad d_T=24d^4,\qquad
 \int_Bw^2=6d^2-\frac1{4d^2}>0.
\]
The integral cycle $\widetilde w=d_T[S]_B-
\kappa^*(t^2-2st+3s^2)$ has
\[
        \int_B\widetilde w^2=144d^6(24d^4-1)>0.
\]
By H9-T7 the corresponding height classes have these exact heights
in its normalization. One may choose $E$ supersingular in every
characteristic. These are positive special-axis classes, not a
classification of nonliftable classes and not a wrong-sign example.

\paragraph{Archive and source location.}
The height interpretation is JAC1964 H9-T7, (H9.18), for cote 44,
Th.~5, scan 004. The primitive qualification is the margin on scan
008. The normal-bundle and self-intersection calculations use Fulton,
\emph{Intersection Theory}, Section 6.3, self-intersection formula,
and the codimension-two blowup normalization already checked in
H9-T7. The existence argument uses the ordinary jet construction
underlying SGA~7~II, XVII, 1.1--1.2 (pp.~213--214), 2.5
(p.~217), 3.2.1--3.2.4 (p.~219), 3.3 (p.~221), 3.7
(p.~234), and 4.2--4.3 (pp.~236--239); its imposed-centre
incidence argument is provided here, not attributed as a separate
statement to SGA. In cote 16 language the calculation verifies $I$
on a specified algebraic primitive subspace. It does not establish
$I(B)$ on its complement or close QH9-axis.
